\subsection{Consolidated evidence view}
Table~\ref{tab:evidence} consolidates the development, external robustness, detector, parity-ablation, and finite-index evidence used by the manuscript.
\begin{table}[!htbp]
\caption{Consolidated evidence stack for the TOMM manuscript.}
\label{tab:evidence}
\small
\begin{tabularx}{\textwidth}{p{2.2cm}p{2.8cm}Y Y}
\toprule
Evidence layer & Scale & Principal outcomes & Manuscript role\\
\midrule
BOSSBase development & 5 seeds; 1,800 calibration trials & 1,800/1,800 recovery; 7,995--8,000 unique covers; development AUC 0.514--0.518 aggregate & Method and operating-point evidence\\
Caltech calibration & 5 seeded resamplings; 1,800 trials & 1,800/1,800 recovery; all 7,000 covers used/seed; max RS 37 & External feasibility and calibrated robustness\\
Caltech holdout & 5 seeded resamplings; 1,200 trials & 1,163/1,200 recovery; seven attacks 150/150; crop 12\% 113/150 & External generalization boundary\\
Final steganalysis & 3 detectors; 5 Caltech resamplings $\times$ 20 splits & Resampling-mean AUC 0.504, 0.513, 0.527 with descriptive ranges & Selection-discrimination characterization\\
RS parity ablation & 5 parity levels; 6,000 holdout trials & 69.08\% at RS0 to 96.92\% at RS128; 970 to 2,676.7 images/session & Reliability--communication tradeoff\\
Finite-index planning & 3 payload sizes at RS128 & Conditional $P_{\mathrm{block}}$ from $5.04\times10^{-43}$ to $1.14\times10^{-25}$ & Capacity-risk quantification\\
\bottomrule
\end{tabularx}
\end{table}

\section{Discussion}
\subsection{What the external evidence establishes}
A single frozen protocol transfers across five seeded Caltech train--index resamplings with fixed $K$, group size, parity, projection-ranking rule, mapping, and detector settings. Moderate Gaussian noise, blur, 8\% crop, and 9-degree rotation remain fully recoverable at the authenticated-message level. The 12\% crop condition reaches the 64-byte RS correction radius and concentrates the 37 message failures. Reed--Solomon coding therefore contributes directly to observed robustness: crop 8\% requires as many as 34 corrected byte positions, while the strongest crop reaches the decoder limit.

The group-bank mechanism also demonstrates the value of complementary cover failures. A five-cover group can remain majority-correct when individual members fail on different transformations. This enables every final external schedule to exercise all 7,000 covers while attaining perfect calibration recovery. The result supports channel qualification at the communication-unit level and preserves broad index participation.

\subsection{Reliability, communication cost, and finite-index planning}
The parity sweep exposes a pronounced reliability--traffic tradeoff. Moving from RS0 to RS128 raises authenticated holdout recovery from 69.08\% to 96.92\%, while mean session traffic rises from 970 to 2,676.7 images. At RS128, the present five-cover signaling construction uses 2,320--3,070 images for 8--64-byte plaintexts, corresponding to 0.0276--0.1668 useful payload bits per transmitted image. These measurements identify communication efficiency as the principal deployment bottleneck of the current operating point. They also reveal an observable session-volume side channel because raw cover count varies with payload length. Fixed-volume batching, payload-length padding, dummy-cover insertion, denser group signaling, adaptive parity allocation, and timing shaping are direct design paths for improving deployment efficiency and reducing length-class exposure. The deterministic inequality $d_k\le a_k$ remains the session-specific feasibility test, while the multinomial model supports pre-session capacity planning under the stated uniform-demand assumption.

\subsection{Detectability as a separate security axis}
The detector analysis separates effect magnitude from resampling stability. SRM-lite produces the smallest mean discrimination above chance ($\Delta\mathrm{AUC}=0.0040$), GLCM increases the separation to 0.0128, and the CNN reaches 0.0266. Across the five seeded train--index resamplings, the CNN means remain confined to 0.5232--0.5301. Because these runs reuse the same underlying Caltech corpus, the cross-seed spread is interpreted as sensitivity to assignment rather than independent external inference.

This ordering is consistent with a distributional selection effect that becomes more visible to learned features than to fixed residual statistics. The image-content result is complemented by the session-volume analysis: cover count varies with payload size, so content-level discrimination and traffic-level observability represent distinct operational signals.

\subsection{Why the cyclic session shift matters}
Finite protected messages can produce nonuniform coded-symbol frequencies. A fixed symbol-to-cluster permutation would translate message-level skew into persistent visual-cluster skew. In ARCIS, the additive session shift makes every fixed coded symbol visit every visual cluster exactly once over each complete $K$-session cycle. The Gray-indexed bijection determines symbol order inside the implementation; the balancing invariant itself follows from cyclic addition of the authenticated session index. This symbol-wise invariant complements the empirical signature-balanced group scheduler. Aggregate traffic balance across several sessions additionally depends on the coded-symbol histograms carried by those sessions.

\subsection{Cross-corpus protocol validation}
The BOSSBase and Caltech results provide cross-corpus evidence at the protocol-design level. BOSSBase supports operating-point selection. On Caltech, the same protocol design is retained while the descriptor, codebook, projection, and group bank are fitted within each Caltech resampling. Calibration then reaches 1,800/1,800 authenticated recoveries while exercising every cover in each 7,000-image index; held-out attack strengths reach 96.92\% overall, with the 12\% crop defining the strongest observed boundary. The experiment therefore supports protocol-level transfer with target-corpus fitting rather than frozen-model transfer.

\subsection{Relation to recent coverless and generative work}
\method and Guo--Ping address closely related natural-image selection problems with different recovery events. Guo--Ping report robustness at 10-, 14-, and 15-bit nominal-capacity settings under their published datasets and native evaluation metric~\cite{guo2026capacity}. ARCIS evaluates 8-, 32-, and 64-byte plaintexts and counts success only when the complete plaintext is recovered after majority decoding, RS correction, framing, and AES-GCM verification. Table~\ref{tab:positioning} therefore separates quantitative regime from recovery semantics instead of combining heterogeneous percentages into one performance scale. DiffStega represents a distinct generative regime in which a secret image is reconstructed with diffusion models~\cite{yang2024diffstega}; its reconstruction evidence remains useful for architectural positioning within that native information object.

\section{Limitations and Threats to Validity}
\textbf{External validity.} Caltech-101 provides one external natural-image corpus, while BOSSBase supports development. The five Caltech runs are seeded train--index resamplings of that same corpus and overlap substantially across seeds, so they measure sensitivity to data assignment rather than independent external replication. Validation on additional native-resolution corpora, social-media recompression pipelines, and heterogeneous acquisition devices would broaden domain evidence.

\textbf{Channel boundary.} The robustness claim covers the declared transformation families and strengths. The 12\% crop condition defines the strongest observed boundary of the present protocol. Perspective warps, larger crops, occlusion, color transformations, arbitrary resampling chains, and learned image editing provide natural directions for broader evaluation.

\textbf{Warden and traffic scope.} The SRM-lite, GLCM, and CNN battery spans residual, texture, and learned selected-vs.-remainder discrimination. The measured transmission counts additionally expose session-volume observability. Semantic side information, temporal correlation, network metadata, and adaptive deployment-specific wardens provide complementary directions for extending this evaluation.

\textbf{Cryptographic scope.} AES-GCM, HMAC, and Reed--Solomon coding are standard components integrated with the finite-index signaling protocol, session balancing, and feasibility construction. The public benchmark key/nonce derivation supports exact reproducibility; operational deployment can substitute institutionally managed keys and nonces.

\textbf{Comparator scope.} DiffStega is retained as executable positioning evidence under secret-image reconstruction semantics, while Guo--Ping is positioned under its native cover-selection capacity and robustness semantics. The recovery-event definitions in Table~\ref{tab:positioning} preserve these distinctions explicitly.

\section{Reproducibility and Artifact Traceability}
The project artifact records dataset provenance, frozen seeds, calibration/holdout protocols, group-bank construction, exact session mapping, detector schedules, workflow identifiers, commands, output manifests, and machine-readable manuscript-facing results. The \path{tomm_results} directory contains \path{rs_parity_ablation_summary.json}, \path{detector_resampling_summary.json}, \path{caltech_resampling_overlap.json}, \path{communication_cost.json}, and \path{finite_index_blocking_risk.json}. The public repository is \url{https://github.com/EkodeckStephane/NARCIS}; the historical repository name is retained for provenance while ARCIS is the current scientific identity.

\section{Conclusion}
\method formulates selection-based coverless image communication as a finite-index authenticated signaling problem. External calibration recovers 1,800/1,800 authenticated messages and the 1,200 unseen holdout trials yield 1,163 recoveries (96.92\%), with all 37 failures concentrated in the 12\% crop. The controlled parity ablation moves from 69.08\% recovery at RS0 to 96.92\% at RS128 while mean transmitted images per session rise from 970 to 2,676.7; at RS128, 8--64-byte plaintexts require 2,320--3,070 images. The deterministic feasibility condition is supplemented by an explicitly conditional blocking-risk model, while resampling-level detector means of 0.504, 0.513, and 0.527 characterize low-amplitude selection discrimination under the evaluated wardens.

These results establish a reproducible operating point with bounded channel robustness, an explicit reliability--communication tradeoff, deterministic finite-index feasibility, conditional blocking-risk quantification, complete-index participation, authenticated recovery, and detector-scoped security evidence. The measured carrier volume identifies communication efficiency and traffic shaping as the principal targets for the next ARCIS operating point.

\begin{acks}
This work received no external financial support. The authors thank their respective institutions for research support.
\end{acks}

\bibliographystyle{ACM-Reference-Format}
\bibliography{ARCIS_references}

\end{document}
